% Fate's Edge SRD −−− GM Guidelines (Section 22)

\section{GM Guidelines}

This chapter consolidates advice for running Fate's Edge smoothly, fairly, and with narrative focus. It covers session preparation, in−play adjudication, common pitfalls, and tools for managing cognitive load.

\subsection{Session Zero}

Before the campaign begins, establish the following with your players.

⫕section{Safety Tools}
\begin{itemize}
  ℹtem \textbf{Lines:} Content that will not appear (e.g., no sexual violence, no harm to children).
  ℹtem \textbf{Veils:} Content that happens off screen or fades to black.
  ℹtem \textbf{X−Card:} Any player may stop a scene at any time without explanation.
\end{itemize}
Record Lines and Veils where everyone can see them. Keep an X−Card visible on the table.

⫕section{Tone and Expectations}
\begin{itemize}
  ℹtem What is the campaign's primary genre? (intrigue, exploration, mercenary, horror, etc.)
  ℹtem How lethal and how grim should the tone be?
  ℹtem How much player vs. GM authority over the story?
  ℹtem Expected session length and attendance.
\end{itemize}

⫕section{Campaign Frame (Optional)}
Consider using the Crown Spread (Section~\ref{sec:crown−spread}) to generate a shared campaign frame. Draw 5–7 cards to establish motifs, factions, omens, and a looming climax.

\subsection{The 30−Second Adjudication Loop}

When a player declares an action, resolve it with this loop:

\begin{enumerate}
  ℹtem \textbf{Clarify intent and approach.} ``What do you want, and how?''
  ℹtem \textbf{Set stakes and Position.} ``If it works, what changes? If it fails, what bites?'' Start Dominant/Standard unless fiction says otherwise.
  ℹtem \textbf{Roll and read.} Count successes (6–9 = 1, 10 = 2); count 1s as SB.
  ℹtem \textbf{Spend one beat well.} Cash SB on one memorable twist or tick a relevant Timer .
  ℹtem \textbf{Push forward.} Describe how the fiction is now different. Ask, ``Who moves next?''
\end{enumerate}

\subsection{Defaults That Keep Things Moving}

When in doubt, use these defaults:

\begin{itemize}
  ℹtem \textbf{Range:} Near (unless otherwise stated).
  ℹtem \textbf{Position:} Controlled.
  ℹtem \textbf{Effect:} Standard.
  ℹtem \textbf{DV:} 3 (2 for trivial, 4 for hard, 5+ for extreme).
  ℹtem \textbf{SB budget per scene:} 4 + highest character Tier. Unspent SB may carry over (max 4).
\end{itemize}

\subsection{When to Roll (and When Not To)}

\begin{itemize}
  ℹtem \textbf{Roll when:} uncertainty exists \textbf{and} meaningful stakes are present \textbf{now}.
  ℹtem \textbf{Do not roll when:}
  \begin{itemize}
    ℹtem The outcome is certain (say ``Yes'').
    ℹtem The action has no real cost or consequence.
    ℹtem The player is fishing for Boons with repetitive, low−stakes rolls.
  \end{itemize}
  ℹtem \textbf{Offer a choice or cost instead of a roll when:} the action is plausible but the player wants to avoid risk. ``You can do it, but it will take time (tick a timer)'' or ``You can do it, but you will mark 1 Fatigue.''
\end{itemize}

\subsection{Timers: The Three−Timer Rule}

At any given time, maintain at most \textbf{three active timers} in a scene:
\begin{enumerate}
  ℹtem \textbf{Situation Timer :} immediate goal (persuasion, escape, research).
  ℹtem \textbf{Counter−Clock:} opposition pressure (guards arriving, ritual completing).
  ℹtem \textbf{Campaign Timer :} long−term front (faction rise, curse spreading).
\end{enumerate}
Merge or retire additional timers. Timers are visible; name them and tick them in front of players.

\subsection{Spending Story Beats (SB)}

\begin{itemize}
  ℹtem \textbf{One strong beat is better than three tiny penalties.} Change the fiction, don't just add numbers.
  ℹtem \textbf{Tie SB spends to visible outcomes:} a new foe appears, a path closes, a timer advances.
  ℹtem \textbf{Surprise yourself.} If you don't know how a complication will resolve, spend SB and find out together.
\end{itemize}

\subsection{Balancing Encounters by Tier}

Fate's Edge does not use Challenge Ratings. Instead, consider the party's Tier and the enemy's Cap.

\begin{tabular}{@{}lll@{}}
⊤rule
\textbf{Party Tier} & \textbf{Typical Threat Cap} & \textbf{Notes} \\
∣rule
Tier I (0–40 XP) & Cap 1–2 & Single enemies; avoid area effects. \\
Tier II (41–90 XP) & Cap 2–3 & Lieutenants with minions. \\
Tier III (91–150 XP) & Cap 3–4 & Bosses with phase timers. \\
Tier IV (151–220 XP) & Cap 4–5 & Mythic threats; require preparation and assets. \\
⊥tomrule
\end{tabular}

⫕section{Building an Enemy}
\begin{itemize}
  ℹtem \textbf{Cap:} Determines dice pools, Harm output, and resistance.
  ℹtem \textbf{Harm:} Typical Harm output = Cap (e.g., Cap 2 deals Harm 2 on a hit).
  ℹtem \textbf{Defense DV:} 3 for standard foes, 4 for elites, 5+ for bosses.
  ℹtem \textbf{Timers:} Bosses should have phase timers (6–10 segments) that change their behavior when filled.
\end{itemize}

\subsection{The Four Questions (When Stuck)}

Ask out loud:

\begin{enumerate}
  ℹtem If this goes right, what changes? (Intent)
  ℹtem If this goes wrong, what bites back? (Stakes)
  ℹtem What single twist will make this memorable? (SB spend)
  ℹtem Who moves next? (Momentum)
\end{enumerate}

\subsection{Common Pitfalls and Fixes}

\begin{tabular}{@{}lp{8cm}@{}}
⊤rule
\textbf{Pitfall} & \textbf{Fix} \\
∣rule
Over−cranking SB (scenes feel punitive) & Halve your SB spends for a while; cash into visible timers instead of immediate penalties. \\
Timer sprawl & Merge redundant timers. Each active scene rarely needs more than 2–3. \\
Tag paralysis (player waits for perfect tag) & Paraphrase: ``Sounds like [VEIL]. DV 3. Want to roll?'' \\
Rules drift & If table memory conflicts with text, pick the ruling that keeps flow, then note to reconcile after play. \\
Player analysis paralysis & Offer two clear options and one risky third. ``You can A, B, or try something bold.'' \\
⊥tomrule
\end{tabular}

\subsection{Handling Player Creativity}

\begin{itemize}
  ℹtem \textbf{Say ``Yes, but'' not ``No.''} Offer a cost, a timer, or a complication.
  ℹtem \textbf{Use Setup and Assist actions.} Let non−specialists contribute meaningfully.
  ℹtem \textbf{Embrace unexpected solutions.} If a player bypasses a planned encounter, let it work — then escalate somewhere else.
  ℹtem \textbf{Fail forward.} A missed roll should never stop the story; it should change direction.
\end{itemize}

\subsection{Session Checklist}

⫕section{Before Play}
\begin{itemize}
  ℹtem Review previous session's fallout and unresolved timers.
  ℹtem Update faction and front timers.
  ℹtem Prepare 2–3 flexible scenes (not scripts).
  ℹtem Set the initial Position and DV for the opening scene.
  ℹtem Remind players of safety tools.
\end{itemize}

⫕section{During Play}
\begin{itemize}
  ℹtem Use the 30−second loop.
  ℹtem Spend one strong beat per roll (not many small penalties).
  ℹtem Keep timers visible and tick them openly.
  ℹtem Rotate spotlight among players.
  ℹtem When you don't know a rule, make a ruling and note it for later.
\end{itemize}

⫕section{After Play}
\begin{itemize}
  ℹtem Award XP (Section~\ref{sec:advancement}).
  ℹtem Advance or clear active timers.
  ℹtem Note rulings to reconcile with rules text.
  ℹtem Ask players: one thing they loved, one thing they want more of.
\end{itemize}

\subsection{Rookie GM Comfort Dials}

If you are new to Fate's Edge, start with these training wheels:

\begin{itemize}
  ℹtem \textbf{Soft SB:} For your first 2 sessions, cap each roll's SB spend to 1–2 unless it is a set−piece.
  ℹtem \textbf{Visible Timers:} Put all timers on the table. Name them aloud. Tick them in ink.
  ℹtem \textbf{Tag Cards:} Print one−liners for frequently used tags ([WARD], [BANISH], [COUNTER]). Hand them out when a power is active.
  ℹtem \textbf{One Move, One Sentence:} Every ruling should end with one sentence that states the new situation.
\end{itemize}

\subsection{Design Guardrails (for Consistency)}

These principles guide rulings and homebrew:

\begin{itemize}
  ℹtem \textbf{Narrative Primacy:} Mechanics serve story, not replace it.
  ℹtem \textbf{Risk as Drama:} Every roll carries potential for triumph and complication.
  ℹtem \textbf{Meaningful Growth:} XP changes characters and the world.
  ℹtem \textbf{Consequence Weight:} Choices ripple outward; nothing is free.
  ℹtem \textbf{Fail Forward:} Misses fuel Boons; 1s become SB.
\end{itemize}

\subsection{Final Advice}

The rules are rails you lay just ahead of the train. Do not let perfect knowledge of the system slow the table. When in doubt, follow the fiction, set a Position and DV, and roll. The game is designed to carry you.

